\begin{center}
\footnotesize
\begin{tabularx}{\textwidth}{>{\raggedright\arraybackslash}X>{\raggedright\arraybackslash}X>{\raggedright\arraybackslash}X>{\raggedright\arraybackslash}X}
\toprule
Result & Basis & Depends on & Reliability \\
\midrule
The rollover clock and the consolidated stock & Security-level data, exact against official totals; backtest error 0.05 pp; the Fed's 2022–25 loss reproduced & Bucket correction for 1980–2002 & High \\
The fiscal threshold does not depend on maturity & Proof (Section 3.2) & Linear debt dynamics; constant $\psi$ & High, within the model class \\
Funding central-bank bonds with reserves raises the price of the inflation route & Accounting; the same in the United States, the UK and Japan; Japan's 2024 reform & Ignores the return on the Fed's MBS & High \\
QE lowered the monetary tolerance threshold $\kappa^*$ & Same accounting, three countries & First order; the response lasts as long as the surprise & High for the fall (about two fifths at H = 5, 10, 15); below 0.5 at H = 10 and 15, not at H = 5 \\
Markets priced far less of the post-2020 surprises than full Fisher repricing & U.S. and UK tests (1.8 times in both); Japan & UK and Japan: closed portfolio & High \\
2023–25 has the largest requirement since 1980 & Year-by-year accounting & Today's $\psi$ at least about 2bp (3bp at $\hat\varphi$ = 0.25) & Medium \\
Levels of the requirement (1.0–1.1 pp; about 4 pp to hold debt/GDP) & Same & $\psi$, g, H; currency demand held fixed & Order of magnitude \\
The level of the fiscal response $\hat\varphi$ & UK fiscal events about 0.15; U.S. not identified & — & Low; not needed for the results \\
\bottomrule
\end{tabularx}
\end{center}
